\documentclass[11pt,a4paper]{article}

\usepackage[a4paper,margin=1in]{geometry}
\usepackage[T1]{fontenc}
\usepackage[utf8]{inputenc}
\usepackage{array}
\usepackage{booktabs}
\usepackage{enumitem}
\usepackage{hyperref}
\usepackage{longtable}
\usepackage{tabularx}
\usepackage{xcolor}
\usepackage{listings}

\hypersetup{
  colorlinks=true,
  linkcolor=blue,
  urlcolor=blue,
  pdftitle={Dell Presales Academy Global Documentation},
  pdfauthor={TCO Repatriation Dashboard Project}
}

\lstset{
  basicstyle=\ttfamily\small,
  breaklines=true,
  columns=fullflexible,
  frame=single,
  rulecolor=\color{black!20},
  backgroundcolor=\color{black!3}
}

\setlist[itemize]{leftmargin=*,itemsep=0.2em}
\setlist[enumerate]{leftmargin=*,itemsep=0.2em}
\renewcommand{\arraystretch}{1.18}

\title{Dell Presales Academy\\Global Documentation Package}
\author{Generated from the current Dell Presales Academy documentation}
\date{May 30, 2026}

\begin{document}

\maketitle
\tableofcontents
\newpage

\section{Document Purpose}

This document consolidates the Dell Presales Academy documentation package for the TCO Repatriation Dashboard into one LaTeX-ready presentation and interview reference. It combines the package README, one-page brief, role alignment, Dell alignment notes, talk tracks, presentation deck, live demo runbook, and Q\&A preparation.

This is not official Dell Technologies documentation. It is a candidate-built package designed to present the TCO Repatriation Dashboard as evidence of presales reasoning, business-value framing, technical architecture, and customer-facing communication.

\section{Package Overview}

\subsection{Audience}

The intended audience includes:

\begin{itemize}
  \item Dell Presales Academy recruiters.
  \item Interviewers for an Inside Product Solution Architect role.
  \item Technical or business reviewers who want to see presales thinking, not only coding ability.
\end{itemize}

\subsection{Core Positioning}

\begin{lstlisting}
This is a candidate-built presales tool that turns infrastructure cost pressure into a structured customer conversation about workload placement, hybrid strategy, and business value.
\end{lstlisting}

\subsection{Source Files}

\begin{longtable}{p{0.34\textwidth}p{0.56\textwidth}}
\toprule
\textbf{Source File} & \textbf{Purpose} \\
\midrule
\texttt{README.md} & Package overview, audience, file map, and project boundary. \\
\texttt{one-page-brief.md} & Recruiter-friendly one-page project summary. \\
\texttt{role-alignment.md} & Mapping between the project and the Inside Product Solution Architect role. \\
\texttt{demo-runbook.md} & Live demo flow for a 7 to 10 minute recruiter presentation. \\
\texttt{qa-prep.md} & Likely interview questions and strong answers. \\
\texttt{dell-alignment.md} & Dell public-message alignment and source notes. \\
\texttt{talk-track.md} & Short, medium, long, Dell-focused, and closing verbal pitches. \\
\texttt{presentation-deck.md} & Slide-by-slide presentation with speaker notes. \\
\bottomrule
\end{longtable}

\subsection{Recommended Use Order}

\begin{enumerate}
  \item Read the one-page brief first.
  \item Practice the 90-second pitch.
  \item Use the slide deck if asked to present formally.
  \item Use the demo runbook when showing the live app.
  \item Review the Q\&A preparation before interviews.
\end{enumerate}

\section{One-Page Brief}

\subsection{Project}

\textbf{TCO Repatriation Dashboard}

\subsection{Positioning}

A presales decision-support dashboard that helps compare public cloud costs against local infrastructure subscription costs and recommends whether a workload is better suited for cloud-first, repatriation, or hybrid placement.

\subsection{Why It Was Built}

The project was built to demonstrate thinking beyond code. It connects technical infrastructure choices with financial outcomes, which is central to presales and solution architecture work.

\subsection{Business Problem}

Organizations often move to public cloud for agility, but some steady-state workloads become expensive over time. Customers need help understanding whether to keep workloads in public cloud, move predictable parts to local infrastructure, or adopt a hybrid strategy.

\subsection{Solution}

The app allows a presales user to enter existing cloud costs and local infrastructure assumptions. It calculates multi-year TCO, break-even timing, savings, and a hybrid placement recommendation.

\subsection{Core Features}

\begin{itemize}
  \item Public cloud cost inputs.
  \item Local infrastructure cost inputs.
  \item 12 to 60 month TCO projection.
  \item Break-even analysis.
  \item Cumulative TCO visualization.
  \item Monthly and cumulative delta charts.
  \item Cost composition chart.
  \item Rule-based hybrid recommendation.
  \item CSV export.
  \item Tests and full documentation.
\end{itemize}

\subsection{Presales Value}

\begin{longtable}{p{0.32\textwidth}p{0.58\textwidth}}
\toprule
\textbf{Presales Skill} & \textbf{How The Project Shows It} \\
\midrule
Customer discovery & Cost categories map to discovery questions. \\
Business value & Savings, payback, and TCO are front and center. \\
Technical architecture & Separates cloud, local, and hybrid placement logic. \\
Executive communication & Charts and summaries are built for non-developers. \\
Trust & Formulas and assumptions are visible. \\
Pragmatism & Recommendation avoids overclaiming exact migration percentages. \\
\bottomrule
\end{longtable}

\subsection{Dell-Relevant Framing}

The project aligns with public Dell themes around practical progress through technology, simplified IT consumption, hybrid and multicloud infrastructure choices, modern data center and edge readiness, and cost-aware enterprise AI infrastructure decisions.

\subsection{Strong Interview Statement}

\begin{lstlisting}
I did not build this as a generic calculator. I built it as a presales conversation tool. The output helps a customer understand the financial impact, the architectural tradeoffs, and the next step for workload placement.
\end{lstlisting}

\section{Talk Tracks}

\subsection{20-Second Version}

\begin{lstlisting}
I built a TCO Repatriation Dashboard to show how I think as a presales solution architect. It compares public cloud and local infrastructure costs, visualizes break-even and savings, and gives a transparent hybrid placement recommendation.
\end{lstlisting}

\subsection{45-Second Version}

\begin{lstlisting}
This project is a presales decision-support tool. The business problem is that customers often realize some public cloud workloads become expensive over time. The dashboard lets a user enter cloud and local infrastructure assumptions, compare multi-year TCO, identify break-even, and get a conservative recommendation: cloud-first, hybrid, repatriation candidate, or validate assumptions. I built it to show that I can connect technical architecture to financial value and customer communication.
\end{lstlisting}

\subsection{90-Second Version}

\begin{lstlisting}
I built this project for the kind of conversation a presales solution architect has with customers. The problem is not simply that cloud is expensive. The real question is workload placement: what should stay in public cloud, what might run better locally, and where does a hybrid strategy make sense?

The dashboard models public cloud costs against local infrastructure subscription costs over a selected projection window. It shows TCO, break-even, net savings, monthly run-rate delta, and cost composition. I also added a rule-based hybrid recommendation that explains why it recommends cloud-first, hybrid, repatriation, or more validation.

I kept the logic transparent because presales recommendations need trust. The formulas are documented, the assumptions are visible, the recommendation is conservative, and the calculation and recommender modules are tested. My goal was to show that I can build something technical while also thinking like a customer-facing architect: business problem first, technical tradeoffs second, recommendation third.
\end{lstlisting}

\subsection{Dell-Focused Version}

\begin{lstlisting}
I aligned the project with the kind of infrastructure conversations Dell is positioned to lead: hybrid and multicloud, private cloud, modern data centers, edge, consumption models, and AI infrastructure economics. I am not presenting it as a Dell product. I am presenting it as evidence that I understand how to connect those themes to a customer-facing business conversation.
\end{lstlisting}

\subsection{Closing Line}

\begin{lstlisting}
The project shows the bridge I want to bring to the role: technical depth, business value, and customer-facing clarity.
\end{lstlisting}

\section{Role Alignment}

\subsection{Role Fit Thesis}

An Inside Product Solution Architect needs to connect customer problems, technical solutions, business value, and sales execution. This project demonstrates that full chain.

\subsection{Capability Mapping}

\begin{longtable}{p{0.40\textwidth}p{0.50\textwidth}}
\toprule
\textbf{Role Capability} & \textbf{Evidence From Project} \\
\midrule
Understand customer business pain & Frames cloud cost pressure as a TCO and workload placement problem. \\
Translate pain into solution criteria & Converts cost categories into cloud and local comparison inputs. \\
Communicate financial impact & Shows TCO, savings, savings percent, run-rate delta, and break-even. \\
Explain technical tradeoffs & Uses hybrid recommendation instead of binary cloud versus on-prem thinking. \\
Support sales conversations & Includes one-page summary, demo script, report notes, and CSV export. \\
Build trust with stakeholders & Keeps assumptions and formulas visible. \\
Handle ambiguity & Uses confidence levels and validation language. \\
Work cross-functionally & Project artifacts support technical, financial, and sales audiences. \\
\bottomrule
\end{longtable}

\subsection{Role Connection Statement}

\begin{lstlisting}
The role is not only about knowing infrastructure products. It is about helping customers make better decisions. I built this project to show that I can take a customer concern, model the business impact, explain the technical options, and guide the next step in a way that sales, finance, and IT stakeholders can all understand.
\end{lstlisting}

\subsection{Presales Behaviors Demonstrated}

\subsubsection{Discovery}

The dashboard input categories mirror discovery questions:

\begin{itemize}
  \item Where is cloud spend concentrated?
  \item How fast is the workload growing?
  \item What discounts already exist?
  \item What would local infrastructure cost to operate?
  \item What one-time migration costs would affect payback?
\end{itemize}

\subsubsection{Qualification}

The recommendation qualifies whether the opportunity is strong enough for repatriation, better positioned as hybrid, not economically justified, or too close to call without better assumptions.

\subsubsection{Value Framing}

The app explains value through total cost, savings, payback timing, cost composition, and risk of late break-even.

\subsubsection{Technical Credibility}

The implementation separates UI, cost models, TCO calculator, recommendation rules, export logic, and tests. This makes the project defensible in a technical review.

\subsection{Interview Positioning}

Avoid saying:

\begin{lstlisting}
I built a Streamlit app.
\end{lstlisting}

Say:

\begin{lstlisting}
I built a presales decision-support tool that uses a technical cost model to guide customer conversations about hybrid workload placement.
\end{lstlisting}

\subsection{What This Says About The Candidate}

The project is designed to show commercial thinking, clear communication, working technical asset creation, infrastructure economics understanding, disciplined recommendation language, and alignment between technical recommendations and customer outcomes.

\section{Dell Alignment Notes}

\subsection{Purpose}

This section explains how the TCO Repatriation Dashboard can be framed for Dell Presales Academy recruiters. It does not claim the project is an official Dell Technologies product.

\subsection{Public Dell Themes Used}

The project was aligned against public Dell and Dell Technologies messaging available as of May 30, 2026.

\subsubsection{Technology Should Drive Practical Progress}

Dell describes its purpose around driving human progress through technology. The dashboard aligns because it helps make infrastructure decisions more understandable, turns technical cost data into a business decision, and supports better customer outcomes through transparent analysis.

Recruiter framing:

\begin{lstlisting}
I tried to build a tool that uses technology to make a customer decision clearer, not just a tool that shows technical features.
\end{lstlisting}

\subsubsection{Simplified Technology Consumption}

Dell APEX messaging emphasizes simplified consumption and as-a-Service models across infrastructure. The project reflects this by modeling local infrastructure as a monthly subscription, using OPEX-style thinking instead of only hardware purchase cost, and helping customers understand predictable consumption over time.

Recruiter framing:

\begin{lstlisting}
The model reflects the way many customers want to consume infrastructure today: with predictable cost, subscription thinking, and simplified decision-making.
\end{lstlisting}

\subsubsection{Hybrid, Multicloud, Data Center, And Edge Reality}

Dell public messaging discusses modern workloads across on-premises data centers, cloud, and edge environments. The recommender avoids binary thinking, frames the decision as workload placement, and supports hybrid recommendations when economics are mixed.

Recruiter framing:

\begin{lstlisting}
I designed the recommendation to be hybrid by default because that is how real infrastructure decisions usually work.
\end{lstlisting}

\subsubsection{Enterprise AI Infrastructure Economics}

Dell AI infrastructure messaging highlights that cost, security, and data challenges can slow AI progress. The same TCO discipline can be applied to AI and data-heavy workloads because cost, data gravity, and security make infrastructure placement more important.

Recruiter framing:

\begin{lstlisting}
AI makes infrastructure placement even more important because cost, data gravity, and security become harder to ignore.
\end{lstlisting}

\subsection{Interview Alignment Statement}

\begin{lstlisting}
I aligned the project with Dell's public direction around hybrid infrastructure, simplified consumption, modern data centers, edge, and AI infrastructure economics. The dashboard is not a Dell product, but it demonstrates that I understand the customer conversations Dell is positioned to lead.
\end{lstlisting}

\subsection{Official Source Notes}

\begin{itemize}
  \item Dell Who We Are page: \url{https://www.dell.com/en-us/lp/dt/who-we-are}
  \item Dell Technologies APEX announcement: \url{https://investors.delltechnologies.com/news-releases/news-release-details/dell-technologies-apex-transforms-how-world-consumes-technology}
  \item Dell Technologies modern data center announcement: \url{https://investors.delltechnologies.com/news-releases/news-release-details/dell-technologies-transforms-data-center-operations-software}
  \item Dell Technologies AI infrastructure announcement: \url{https://investors.delltechnologies.com/news-releases/news-release-details/dell-technologies-fuels-enterprise-ai-innovation-infrastructure}
\end{itemize}

\subsection{Boundaries}

Do not say:

\begin{lstlisting}
This is a Dell solution.
\end{lstlisting}

Say:

\begin{lstlisting}
This is a candidate project aligned to the types of customer problems Dell helps solve.
\end{lstlisting}

Do not say:

\begin{lstlisting}
This proves Dell is cheaper than cloud.
\end{lstlisting}

Say:

\begin{lstlisting}
This helps structure a workload placement and TCO discussion.
\end{lstlisting}

\section{Presentation Deck}

\subsection{Deck Context}

\begin{longtable}{p{0.30\textwidth}p{0.60\textwidth}}
\toprule
\textbf{Item} & \textbf{Value} \\
\midrule
Title & TCO Repatriation Dashboard. \\
Subtitle & A presales decision-support tool for cloud, local infrastructure, and hybrid workload placement. \\
Target audience & Dell Presales Academy recruiters and Inside Product Solution Architect interviewers. \\
Recommended length & 8 to 10 minutes, plus Q\&A. \\
\bottomrule
\end{longtable}

\subsection{Slide 1: Title}

\textbf{TCO Repatriation Dashboard}. A presales project that connects infrastructure architecture to business value.

Speaker note:

\begin{lstlisting}
This project is not just a coding exercise. I built it to demonstrate how I would support a customer conversation: understand the business problem, quantify the impact, explain options, and recommend a practical next step.
\end{lstlisting}

\subsection{Slide 2: The Customer Problem}

\textbf{Cloud-first does not always mean cloud-only.} Many organizations adopted public cloud for speed and flexibility. Over time, predictable workloads can become expensive, especially when storage, backup, compute, support, and data growth compound over multiple years.

Customer question:

\begin{lstlisting}
Should we keep this workload in public cloud, move it local, or use a hybrid approach?
\end{lstlisting}

Speaker note:

\begin{lstlisting}
The business pain is not simply "cloud is expensive." The real issue is workload placement. Some workloads belong in cloud. Some may be better local. Most real customers need a rational hybrid discussion.
\end{lstlisting}

\subsection{Slide 3: Why This Matters In Presales}

\textbf{A presales architect must translate technology into business outcomes.}

This project shows discovery thinking, TCO and ROI framing, technical tradeoff analysis, hybrid placement strategy, and clear communication for technical and financial stakeholders.

Speaker note:

\begin{lstlisting}
An inside product solution architect needs to do more than explain products. They need to qualify the problem, shape the solution, and help the customer understand the financial and operational impact of the decision.
\end{lstlisting}

\subsection{Slide 4: Solution Overview}

The dashboard compares two cost paths: stay in public cloud or use local infrastructure on a subscription model. It calculates multi-year TCO, break-even month, net savings, monthly run-rate delta, and hybrid placement recommendation.

Speaker note:

\begin{lstlisting}
I intentionally kept the model transparent. The customer can see the assumptions, change the inputs, challenge the recommendation, and export the projection.
\end{lstlisting}

\subsection{Slide 5: Demo Architecture}

\begin{lstlisting}
Presales user -> Streamlit dashboard
Streamlit dashboard -> Cloud and local cost inputs
Cloud and local cost inputs -> TCO calculation engine
TCO calculation engine -> Monthly projection and break-even
Monthly projection and break-even -> Hybrid placement recommender
Monthly projection and break-even -> Charts and CSV export
\end{lstlisting}

Speaker note:

\begin{lstlisting}
The code is structured like a real solution: UI, calculation engine, recommender, export, tests, and documentation. I separated the financial logic from the UI so it can be tested and reused.
\end{lstlisting}

\subsection{Slide 6: What The Dashboard Shows}

Core outputs are public cloud TCO, local infrastructure TCO, net savings, break-even month, cumulative TCO chart, cost composition chart, monthly and cumulative delta charts, and CSV export.

Speaker note:

\begin{lstlisting}
The goal is to help a customer understand not just the final number, but the journey. A lower monthly local run-rate may still be a bad decision if the migration cost delays payback too far.
\end{lstlisting}

\subsection{Slide 7: Hybrid Recommendation}

\textbf{The recommendation avoids fake precision.} It does not claim exact migration percentages. It says cloud-first, hybrid candidate, repatriation candidate, or validate assumptions.

Why this matters:

\begin{itemize}
  \item Real customers rarely move everything.
  \item Workload placement depends on cost, risk, operations, data gravity, and elasticity.
  \item A conservative recommendation is more credible than an overconfident one.
\end{itemize}

Speaker note:

\begin{lstlisting}
I built the recommender as a transparent rule-based system. That is important in presales because customers need to trust the logic, not just see a black-box answer.
\end{lstlisting}

\subsection{Slide 8: Dell-Aligned Thinking}

The project aligns with public Dell Technologies themes: human progress through useful technology, simplified technology consumption, hybrid and multicloud workload placement, modern data center and edge readiness, and cost-aware AI and infrastructure decisions.

Speaker note:

\begin{lstlisting}
I am not presenting this as a Dell product. I am showing that I understand the kind of customer conversations Dell is positioned to lead: hybrid infrastructure, consumption models, private cloud, edge, data protection, and AI infrastructure economics.
\end{lstlisting}

\subsection{Slide 9: Presales Skills Demonstrated}

\begin{longtable}{p{0.30\textwidth}p{0.60\textwidth}}
\toprule
\textbf{Skill} & \textbf{Evidence In Project} \\
\midrule
Discovery & Input model separates cloud and local cost drivers. \\
Business value & TCO, savings, payback, and break-even. \\
Architecture & Modular design and hybrid placement logic. \\
Communication & Executive summary, charts, and report notes. \\
Credibility & Transparent formulas and test coverage. \\
Customer empathy & Conservative recommendation and assumption validation. \\
\bottomrule
\end{longtable}

Speaker note:

\begin{lstlisting}
The project is meant to show how I think as a solution architect: technical enough to build the model, commercial enough to explain the value, and careful enough not to oversell the recommendation.
\end{lstlisting}

\subsection{Slide 10: Live Demo Flow}

\begin{enumerate}
  \item Choose a preset.
  \item Show TCO headline metrics.
  \item Explain break-even.
  \item Open charts.
  \item Open the Recommendation tab.
  \item Show assumptions.
  \item Export CSV.
\end{enumerate}

Speaker note:

\begin{lstlisting}
I would keep the live demo short. The purpose is not to click every field. The purpose is to show that the tool supports a customer conversation from cost pain to next action.
\end{lstlisting}

\subsection{Slide 11: What To Improve Next}

If this became a real presales asset, the next improvements would be PDF executive report, scenario comparison, customer-ready discovery checklist, optional billing CSV import, Dell solution mapping layer, and saved scenarios for account teams.

Speaker note:

\begin{lstlisting}
I avoided overbuilding the MVP. The strongest next step would be turning the current output into a polished customer-facing report and adding a structured discovery workflow.
\end{lstlisting}

\subsection{Slide 12: Closing}

\textbf{This project demonstrates the bridge I want to build in the role.}

\begin{lstlisting}
Technical architecture + customer discovery + financial value + clear recommendation
\end{lstlisting}

Closing statement:

\begin{lstlisting}
I built this to show how I would approach presales: understand the customer's business pressure, translate it into a technical and financial model, and guide them toward a practical solution.
\end{lstlisting}

\section{Demo Runbook}

\subsection{Goal}

Deliver a concise 7 to 10 minute demo that makes the project look like a presales asset, not a hobby app.

\subsection{Setup}

Start the app:

\begin{lstlisting}
streamlit run app.py --server.port 8502
\end{lstlisting}

Open:

\begin{lstlisting}
http://localhost:8502
\end{lstlisting}

Recommended preset:

\begin{lstlisting}
Storage-heavy archive
\end{lstlisting}

This preset makes the hybrid conversation clearer because storage-heavy workloads are easier to explain as local evaluation candidates.

\subsection{Demo Flow}

\subsubsection{1. Business Setup: 45 Seconds}

\begin{lstlisting}
The business problem is cloud cost pressure. Customers do not just ask whether cloud is expensive. They ask where each workload should live. This tool helps structure that conversation with a TCO model and a hybrid placement recommendation.
\end{lstlisting}

\subsubsection{2. Show Inputs: 60 Seconds}

Open the sidebar and show cloud cost categories, local infrastructure categories, and projection window.

\begin{lstlisting}
I kept the model assumption-driven. A presales user can enter known monthly cloud spend, local subscription assumptions, growth, discount, and migration costs. The point is to make the assumptions visible.
\end{lstlisting}

\subsubsection{3. Show Overview: 2 Minutes}

Point to public cloud TCO, local infrastructure TCO, net savings, and break-even.

\begin{lstlisting}
This gives the executive view. The key metric is not only savings. Break-even matters because migration cost can make a lower monthly run-rate unattractive if payback is too late.
\end{lstlisting}

Then show cumulative TCO, monthly run-rate delta, cumulative cost advantage, and cost composition.

\subsubsection{4. Show Recommendation: 2 Minutes}

Open the Recommendation tab.

\begin{lstlisting}
I separated this from the overview so it reads as guidance, not just a metric. The recommender is conservative and transparent. It does not pretend to know exact workload percentages. It gives a credible next step: cloud-first, hybrid, repatriation candidate, or validate assumptions.
\end{lstlisting}

Point to recommendation label, confidence, why this result, and what to do first.

\subsubsection{5. Show Assumptions: 60 Seconds}

Open the Assumptions tab.

\begin{lstlisting}
This tab is important for trust. A customer or account team can challenge every input behind the model.
\end{lstlisting}

\subsubsection{6. Show Export: 30 Seconds}

Click \texttt{Download CSV}.

\begin{lstlisting}
The output can be handed to finance or used for follow-up analysis.
\end{lstlisting}

\subsubsection{7. Close: 45 Seconds}

\begin{lstlisting}
The purpose of the project is to show how I approach presales: understand the customer's business pressure, quantify the impact, explain the architecture tradeoffs, and recommend a practical next action.
\end{lstlisting}

\subsection{What Not To Do}

\begin{itemize}
  \item Do not explain every code file in the demo.
  \item Do not oversell the recommender as artificial intelligence.
  \item Do not claim this is a Dell product.
  \item Do not frame repatriation as always better than cloud.
\end{itemize}

\subsection{Strong Closing Line}

\begin{lstlisting}
This project shows the bridge I want to bring to the role: technical architecture, business value, and customer-facing clarity.
\end{lstlisting}

\section{Q\&A Preparation}

\subsection{Why did you build this project?}

\begin{lstlisting}
I wanted to build something that reflects presales work, not just software development. Cloud cost pressure is a real business issue, and the right answer is often not cloud or on-premises, but workload placement. This project lets me show technical modeling, business value framing, and customer-facing recommendation skills.
\end{lstlisting}

\subsection{How does this relate to Dell?}

\begin{lstlisting}
Dell is strongly positioned around hybrid infrastructure, modern data centers, private cloud, edge, as-a-Service consumption, and enterprise AI infrastructure. This project aligns with those conversations because it helps customers reason about where workloads should run and what the financial impact is.
\end{lstlisting}

\subsection{Is this a Dell product?}

\begin{lstlisting}
No. It is a candidate-built project. I use Dell's public messaging only as alignment context. The project demonstrates how I think about customer problems Dell is well positioned to solve.
\end{lstlisting}

\subsection{Why not just compare AWS pricing to Dell hardware pricing?}

\begin{lstlisting}
That would look more precise but would be less credible without real customer workload data and vendor quotes. I designed the tool around customer-provided assumptions so the conversation stays transparent. In a real sales motion, this could be extended with validated pricing and account-specific inputs.
\end{lstlisting}

\subsection{Why did you add a hybrid recommendation?}

\begin{lstlisting}
Because most real infrastructure decisions are not binary. Some workloads need cloud elasticity or managed services. Others are predictable and may be better suited to local infrastructure. The hybrid recommendation makes the tool more realistic and closer to how a solution architect would guide a customer.
\end{lstlisting}

\subsection{Why is the recommender rule-based instead of AI-based?}

\begin{lstlisting}
For a presales tool, explainability matters. A rule-based recommender lets the customer see the logic. It avoids black-box recommendations and prevents false precision.
\end{lstlisting}

\subsection{What would you improve next?}

\begin{lstlisting}
I would add a customer-ready PDF report, scenario comparison, and an optional billing CSV import. I would also add a Dell solution mapping layer if this were used in a Dell-specific context, but I avoided pretending this candidate project has official product pricing or configuration logic.
\end{lstlisting}

\subsection{What technical choices did you make?}

\begin{lstlisting}
I used Streamlit because the goal was fast presales workflow validation. I separated the calculation engine, models, recommender, presets, and exports into modules so the logic is testable and could later support a React frontend or API.
\end{lstlisting}

\subsection{How did you make sure the numbers are reliable?}

\begin{lstlisting}
The formulas are isolated in the calculator module and covered by tests. I added tests for no-growth totals, discount handling, delayed break-even, cases where break-even is not reached, and the main recommender paths.
\end{lstlisting}

\subsection{What does this show about you as a candidate?}

\begin{lstlisting}
It shows I can connect technical architecture to customer value. I can build the tool, explain the business problem, communicate tradeoffs, and avoid overclaiming. That is the kind of discipline I think matters in presales.
\end{lstlisting}

\subsection{Hard Question: Is this too simple?}

\begin{lstlisting}
The calculation is intentionally simple enough to be explainable. In presales, a model the customer trusts is more useful than a complex model they cannot challenge. I focused the MVP on transparency, then documented the path to deeper functionality like billing import, PDF output, and solution mapping.
\end{lstlisting}

\subsection{Hard Question: What if cloud is still cheaper?}

\begin{lstlisting}
Then the recommendation should say cloud-first. The goal is not to force repatriation. The goal is to guide the customer toward the right placement decision. A credible presales architect should be willing to tell the customer when the proposed direction is not justified.
\end{lstlisting}

\subsection{Hard Question: What makes this presales and not just finance?}

\begin{lstlisting}
The tool does not stop at cost. It connects cost to workload placement and architecture choices. It shows when to keep elastic workloads cloud-side, when to evaluate storage-heavy workloads locally, and when to validate assumptions before recommending anything.
\end{lstlisting}

\section{Consolidated Analysis}

\subsection{What The Package Communicates Well}

\begin{itemize}
  \item It frames the project as a presales asset rather than a generic calculator.
  \item It repeatedly ties technical implementation to customer-facing business value.
  \item It avoids overclaiming Dell affiliation, product pricing, or official solution status.
  \item It includes multiple reusable formats: short pitch, long pitch, formal deck, live demo, and Q\&A.
  \item It makes the candidate's reasoning visible: discovery, qualification, value framing, technical credibility, and communication.
\end{itemize}

\subsection{Main Presentation Risks}

\begin{itemize}
  \item The candidate must clearly state that this is not an official Dell product.
  \item The tool should not be framed as proving local infrastructure is always cheaper.
  \item The recommender should not be described as AI.
  \item The demo should not spend too much time on code-level implementation unless asked.
  \item The financial model should be described as a planning estimate, not a final quote.
\end{itemize}

\subsection{Best Use In An Interview}

\begin{enumerate}
  \item Open with the 20-second or 45-second pitch.
  \item Use the one-page brief framing to explain business value.
  \item Demo only the workflow that supports the customer conversation.
  \item Use the role-alignment language to connect the project directly to the Inside Product Solution Architect role.
  \item Use the Q\&A answers to stay credible and avoid overstating the project.
\end{enumerate}

\section{Conclusion}

The Dell Presales Academy package positions the TCO Repatriation Dashboard as a candidate-built presales decision-support tool. Its strongest value is not only the app itself, but the way the documentation frames a customer problem, explains business impact, maps technical work to solution architecture behaviors, and preserves credibility through clear boundaries.

\end{document}
